% Student Number: 2240357068
% Student Name: Baturay KAFKAS
% EEE @ Hacettepe University

% Last update: 09:49 29/10/25

% My resource collection: https://github.com/kfksbtry/uni
% Circuits were set up in LTspice.

\documentclass{article}

\usepackage{graphicx}
\usepackage[top=25mm, bottom=25mm, left=25mm, right=25mm]{geometry}
\usepackage{amsmath}
\usepackage{moresize}
\usepackage{parskip}
\usepackage{float}
\usepackage{fancyhdr}
\usepackage{booktabs}

\pagestyle{fancy}
\fancyhf{}
\fancyhead[R]{Baturay KAFKAS 2240357068 Electrical \& Electronics Engineering}

\rfoot{\thepage}
\renewcommand{\headrulewidth}{0pt} 
\renewcommand{\footrulewidth}{0pt}

\begin{document}

\large
\textit{This part of the experiment is prepared with Online LaTeX Editor Overleaf and the circuits are drawn in LTspice. Visit the website for the source here:}

\textbf{https://www.overleaf.com/read/qgkyrtmrpqkd\#2fdec2}

\hrule

\hfill

\textbf{3. EXPERIMENT 1 - PRELIMINARY WORK}

\textbf{3.1} Calculate the values of the currents $I$, $I_{1}$, $I_{2}$, $I_{3}$ and $I_{4}$ and the voltages $V_{1}$, $V_{2}$, $V_{3}$, $V_{4}$ and $V_{5}$ in \textit{Figure 6}. Show the meter connections for these measurements.

\begin{figure}[H]
    \centering
    \includegraphics[width=0.7\linewidth]{msedge_FufYaRiyNq.png}
\end{figure}

\textbf{Answer}: The connections of meters should be as follows.

\begin{figure}[H]
    \centering
    \includegraphics[width=0.9\linewidth]{1.png}
\end{figure}

\hfill

The ammeters are connected properly so the current enters from the positive side. As for the type of connection, the ammeters are connected in series. Meanwhile, the voltmeters are connected in parallel to the resistors. This will help us accurately measure the desired values.

Let's calculate \textit{I}. Since the main current has not yet been distributed, \textit{I} is equal to the current provided by the source.

\[I = 10 \text{ mA}\]

The equivalent resistance with respect to the terminals $a-b$ is $330\parallel(220+(220\parallel(100+120)))=330\parallel330$. From the current divider, we may calculate $I_1$.

\[I_1=10\cdot10^{-3}\cdot\frac{330}{330+330}=5\text{ mA}\]

From here $I_2=5\text{ mA}$. From the current divider, we may also calculate $I_3$.

\[I_3=5\cdot10^{-3}\cdot\frac{220}{220+100+120}=2.5\text{ mA}\]

Therefore, $I_4=2.5\text{ mA}$.

\hfill

We may now use Ohm's Law to calculate the voltages across the resistors.
\[V_1=330\cdot5\cdot10^{-3}=1.65 \text{ V},\quad V_2=220\cdot5\cdot10^{-3}=1.1 \text{ V},\quad V_3=220\cdot2.5\cdot10^{-3}=0.55\text{ V}\]
\[V_4=100\cdot2.5\cdot10^{-3}=0.25 \text{ V},\quad V_5=120\cdot2.5\cdot10^{-3}=0.3 \text{ V},\]

\[\boxed{\begin{array}{c}I= 10\text{ mA},\quad I_1=5\text { mA},\quad I_2=5\text{ mA},\quad I_3=2.5\text{ mA},\quad I_4=2.5\text{ mA}\\V_1=1.65\text{ V},\quad V_2=1.1\text{ V},\quad V_3=0.55\text{ V},\quad V_4=0.25\text{ V},\quad V_5=0.3\text{ V}\end{array}}\]

{ \textbf{3.2} Find the equivalent resistance $R_{ab}$ between points a and b for the circuit in \textit{Figure 7}.}

\begin{figure}[H]
    \centering
    \includegraphics[width=0.5\linewidth]{msedge_yBjkBAOVeX.png}
\end{figure}

{ \textbf{Answer}: Recall that $R_{eq}$ is equal to the sum of the resistance of the resistors connected in series. $R_{eq}$ is equal to the reciprocal of the sum of the reciprocal of each resistance of the resistors connected in parallel. In mathematical notation:}

\[\displaystyle R_{eq, s} = \sum_{i=1}^{n} R_{i} = R_1 + R_2 + {...} + R_n \quad  \frac{1}{R_{eq, p}} = \sum_{i=1}^{n} \frac{1}{R_{i}} = \frac{1}{R_{1}} + \frac{1}{R_{2}} + {...} + \frac{1}{R_{n}}\]

We approach this question by reducing the number of resistors to obtain an equivalent circuit with respect to the terminals $a$ and $b$.

\begin{figure}[H]
    \centering
    \includegraphics[width=0.5\linewidth]{msedge_Lb39u9VtuS_3.png}
\end{figure}
\[R_{eq, left} = 1 + 2 = 3 \Omega \quad R_{eq, right} = 1 + 2 + 1 = 4\:\Omega\]

\begin{figure}[H]
    \centering
    \includegraphics[width=0.45\linewidth]{yn3Ua6R4B5.png}
\end{figure}

\[\frac{1}{R_{eq}} = \frac{1}{6} + \frac{1}{3} \implies R_{eq} = 2\:\Omega\]

\begin{figure}[H]
    \centering
    \includegraphics[width=0.5\linewidth]{IE5BzPI4P7.png}
\end{figure}

\[R_{eq} = 1 + 2 + 1 = 4\:\Omega\]

\begin{figure}[H]
    \centering
    \includegraphics[width=0.5\linewidth]{NZoVakjMbM.png}
\end{figure}

\[\frac1{R_{eq}} = \frac{1}{4} + \frac{1}{4} \implies R_{eq} = 2\:\Omega\]

\begin{figure}[H]
    \centering
    \includegraphics[width=0.4\linewidth]{J5E4VKKIEB.png}
\end{figure}

\[R_{eq} = 5 + 2 + 5 = 12\:\Omega\]

\begin{figure}[H]
    \centering
    \includegraphics[width=0.4\linewidth]{mEitjTe87Z.png}
\end{figure}

\[\frac1{R_{eq}} = \frac{1}{12} + \frac{1}{12} \implies R_{eq} = 6\:\Omega\]

\begin{figure}[H]
    \centering
    \includegraphics[width=0.14\linewidth]{Draft6.png}
\end{figure}

Therefore, we conclude that $\boxed{R_{ab} = 6\:\Omega}$.

\hfill

{ \textbf{3.3} \textit{Figure 10} shows a carbon resistor. For each of the resistance values given, write down the correct color bands.}

\begin{figure}[H]
    \centering
    \includegraphics[width=0.75\linewidth]{msedge_u1cstyvk4D.png}
\end{figure}

\text{a) 100 $\pm$ 10\% \quad b) 120 $\pm$ 5\% \quad c) 220 $\pm$ 20\% \quad d) 330 $\pm$ 10\%}

\textbf{Answer}: Looking up the color chart, the color bands are given below.

\begin{center}
    \begin{tabular}{ |c|c c c c| }
    \hline
        Option\textbackslash Band& A & B & C & D \\
        \hline
        a)& Brown & Black & Brown & Silver\\
        \hline
        b)& Brown & Red & Brown & Gold\\
        \hline
        c)& Red &Red & Brown & None\\ 
        \hline
        d)& Orange & Orange & Brown & Silver\\ 
        \hline
    \end{tabular}
\end{center}

\hfill

\textbf{3.4} For the circuit given in \textit{Figure 9}, find the Thévenin and Norton equivalent circuits with respect to terminals $a$ and $b$.

\begin{figure}[H]
    \centering
    \includegraphics[width=0.5\linewidth]{msedge_GUhwiHzURX.png}
\end{figure}

\textbf{Answer}: Zero the sources, then we get

\begin{figure}[H]
    \centering
    \includegraphics[width=0.5\linewidth]{2.png}
\end{figure}

$R_{th}=R_{ab}=(1k)\parallel(1k)\parallel(1k)=\dfrac{1000}3\:\Omega$

Analyze the open circuit and find $V_{th}$. Apply the node voltage method by choosing $b$ as the ground.

\[-10\text{ m}+\frac{V_a}{1\text{ k}}+\frac{V_a}{1\text{ k}}+\frac{V_a-12}{1\text{ k}}=0\implies V_a=\frac{22}{3}\text{ V}\]

Therefore, $V_{th}=V_a-0=\dfrac{22}3\text{ V}$. We obtain the Norton-equivalent circuit by $I_{sc}=\dfrac{V_{th}}{R_{th}}=\dfrac{\frac{22}3}{\frac{1000}3}=\dfrac{11}{500}\text{ A}=22 \text{ mA}$. The Thévenin- and Norton-equivalent circuits are as follows, respectively.

\begin{figure}[H]
    \centering
    \includegraphics[width=0.3\linewidth]{3.png} \hspace{30pt}
    \includegraphics[width=0.3\linewidth]{4.png}
\end{figure}

\newpage

\textbf{3.5} Using the superposition principle, find $V_{ab}$ given in \textit{Figure 10}.

\begin{figure}[H]
    \centering
    \includegraphics[width=0.5\linewidth]{msedge_eZ4XQKpVOv.png}
\end{figure}

\textbf{Answer}: First, deactivate $10$ V source, i.e., replace it with a short circuit. Assign a node-voltage label to the top-left node, say $V_c$. Then

\[10\text{ m} + \frac{V_c}{220}+\frac{V_c}{330+100}=0\implies V_c=-1.45\text{ V}\]

The voltage across $100\:\Omega$ is $100\cdot\dfrac{-1.45}{330+100}5=-0.34 \text{ V}$

Deactivate the current source, i.e., replace it with an open circuit. Let $I$ be the circulating current. Applying KVL, we get

\[220I-10+330I+100I=0\implies I=\frac1{65}\text{ A}\]

The voltage across the resistor becomes $100\cdot\dfrac1{65}=1.54 \text{V}$

Using the superposition principle, we sum up the values algebraically. Therefore,

\[\boxed{V_{ab}=-0.34+1.54=1.2\text{ V}}\]

\textbf{3.6} For the circuit given in \textit{Figure 9}, find the dissipated and generated powers.

\textbf{Answer}: We have the Thévenin voltage from $3.4$.

\[I_1=\frac{V_{th}}{1\text{ k}\Omega}=7.33\text{ mA},\qquad I_2=\frac{V_{th}}{1\text{ k}\Omega}=7.33\text{ mA},\qquad I_3=\frac{V_{th}-12}{1\text{ k}\Omega}=-4.67\text{ mA}\]

\[\boxed{P_{10 \text{ mA}}=-I\cdot V_{10\text{ mA}}=10\text{ m}\cdot7.33=-73.3\text{ mW},\quad P_{1000\:\Omega}=\frac{V^2_1}{1000}=\frac{(7.33)^2}{1000}=53.7\text{ mW}}\]
\[\boxed{P_{1000\:\Omega}=\frac{V^2_2}{1000}=\frac{(7.33)^2}{1000}=53.7\text{ mW},\quad P_{1000\:\Omega}=\frac{(V_{th}-12)^2}{1000}=\frac{(-4.67)^2}{1000}=21.8\text{ mW}}\]
\[\boxed{P_{12\text{ V}}=(-4.67 \text{ m})(12)=-56.0\text{ mW}}\]

\textbf{3.7} Show that the total power generated by the sources in \textit{Figure 9} is equal to the total 
power dissipated by the resistors. 

\textbf{Answer}: The small error is due to the approximations based on the power values.

\[P_{dis}=53.7+53.7+21.8=129.2 \text{ mW},\qquad P_{gen}=-73.3-56.0=-129.3 \text{ mW}\]
\[|129.2\text{ mW}|\approx|{-129.3\text{ mW}}|\implies\sum|P_{dis}|=\sum|P_{gen}|\]

\textbf{3.8} Determine the value of $R_L$ that is to be connected between $a$ and $b$ to realize the maximum power transfer.

\textbf{Answer}: Set $R_L = R_{th}$ to obtain maximum power transfer to $R_L$. This comes from the fact that when we write a power function for $R_L$ and take the derivative with respect to $R_L$, we notice that $R_L=R_{th}$ gives rise to a maximum. Therefore, $\boxed{R_L=333.3\:\Omega}$.

\end{document}